\documentclass[12pt,a4paper]{article}

\usepackage[utf8]{inputenc}
\usepackage[T1]{fontenc}
\usepackage{mathptmx} % Times New Roman en texto y simbolos matematicos (guia de catedra)
\usepackage[spanish,es-noquoting]{babel}

\usepackage{amsmath}
\usepackage{amssymb}
\usepackage{graphicx}
\usepackage{float}
\usepackage{booktabs}
\usepackage{geometry}
\usepackage{siunitx}
\usepackage[hidelinks]{hyperref}

\geometry{
    left=2.5cm,
    right=2.5cm,
    top=2.5cm,
    bottom=2.5cm
}

\graphicspath{{../data/figures/}}

\newcommand{\va}{\ensuremath{v_a}}

% ---------------------------------------------------------
% DOCUMENTO
% ---------------------------------------------------------

\begin{document}

\begin{titlepage}

    \centering

    \vspace*{2cm}

    {\Large Instituto Tecnológico de Buenos Aires\par}

    \vspace{1.5cm}

    {\Huge \textbf{Trabajo Práctico 2}\par}

    \vspace{0.5cm}

    {\LARGE Simulación de Sistemas\par}

    \vspace{0.4cm}

    {\large Autómata off-lattice: bandadas de agentes autopropulsados\par}

    \vspace{1.5cm}

    {\large Comisión: G10S2\par}

    \vspace{1.5cm}

    \begin{tabular}{l c}
        \textbf{Integrante} & \textbf{Legajo} \\[0.2cm]
        Javier Liu & 64332 \\
        Juan Ignacio Cantarella & 64509 \\
        Nicolás Rivas & 64292
    \end{tabular}

    \vfill

    {\large 2026 -- Segundo Cuatrimestre\par}

\end{titlepage}

\tableofcontents

\newpage

% =========================================================
\section{Introducción}
% =========================================================

Los sistemas de partículas autopropulsadas presentan un fenómeno colectivo que
no puede deducirse de la regla individual que los gobierna: a partir de
interacciones puramente locales, y sin ninguna dirección privilegiada impuesta
desde afuera, el conjunto puede adoptar espontáneamente una dirección de
movimiento común. Vicsek \textit{et al.} \cite{vicsek1995} mostraron que ese
ordenamiento aparece a través de una transición de fase continua controlada por
un único parámetro de ruido $\eta$, análoga a la transición ferromagnética pero
en un sistema fuera del equilibrio, donde las ``partículas'' que se alinean
además se desplazan y por lo tanto cambian de vecinos a cada paso.

El presente trabajo caracteriza esa transición mediante simulación. Se estudian
dos reglas de interacción distintas sobre la misma dinámica y la misma
geometría: el \textbf{modelo estándar} de Vicsek \cite{vicsek1995}, en el que
cada partícula adopta la dirección promedio de su entorno, y el \textbf{modelo
de votante} \cite{loscar2021}, en el que cada partícula copia la dirección de un
único vecino elegido al azar. La diferencia entre ambos es únicamente cómo se
combina la información de los vecinos --- promediar frente a copiar ---, lo que
permite aislar el efecto de ese mecanismo sobre el ordenamiento colectivo.

Se miden dos observables. El primario es la \textbf{polarización} \va, que
cuantifica cuán alineado está el sistema. El segundo es la \textbf{fracción de
partículas en la componente gigante} $S$ del grafo de vecinos, que cuantifica
cuán conectado está espacialmente. El objetivo es determinar, para las
densidades $\rho = 2,\,4,\,8$, cómo depende cada observable de $\eta$, cómo se
relacionan entre sí y en qué se diferencian las dos reglas de interacción.
Finalmente se evalúa el costo computacional del \textit{Cell Index Method} (CIM)
empleado para la búsqueda de vecinos, contrastándolo con las mediciones del
Trabajo Práctico 1.

% =========================================================
\section{Modelo}
% =========================================================

\subsection{Dinámica}

El sistema consta de $N$ partículas puntuales dentro de una caja cuadrada de
lado $L$ con condiciones periódicas de contorno. Todas se desplazan con una
rapidez común $v$, de modo que el estado de la partícula $i$ queda determinado
por su posición $\mathbf{r}_i(t)$ y por el ángulo $\theta_i(t)$ que forma su
velocidad con el eje $x$. La densidad del sistema es

\begin{equation}
    \rho = \frac{N}{L^2},
    \label{eq:densidad}
\end{equation}

donde $N$ es el número de partículas y $L$ el lado de la caja.

La actualización es \textbf{sincrónica}: las magnitudes en $t+1$ se calculan
todas a partir del estado en $t$. Las posiciones evolucionan según la ecuación
(1) de \cite{vicsek1995},

\begin{equation}
    \mathbf{r}_i(t+1) = \mathbf{r}_i(t) + \mathbf{v}_i(t)\,\Delta t,
    \label{eq:posicion}
\end{equation}

donde $\mathbf{r}_i$ es la posición de la partícula $i$, $\mathbf{v}_i$ su
velocidad --- de módulo $v$ constante y dirección $\theta_i(t)$ --- y $\Delta t$
el paso temporal. Nótese que el desplazamiento se realiza con la dirección
\textbf{en el instante $t$}, no con la recién calculada. Las coordenadas
resultantes se reingresan a la caja por las condiciones periódicas.

Los ángulos evolucionan como

\begin{equation}
    \theta_i(t+1) = \Theta_i(t) + \Delta\theta_i,
    \qquad
    \Delta\theta_i \sim \mathcal{U}\!\left[-\tfrac{\eta}{2},\, \tfrac{\eta}{2}\right],
    \label{eq:angulo}
\end{equation}

donde $\Theta_i(t)$ es la dirección de referencia que fija la regla de
interacción, $\eta$ es la amplitud del ruido y $\Delta\theta_i$ una variable
aleatoria uniforme e independiente para cada partícula y cada paso. El
observable $\eta$ se expresa en radianes.

\subsection{Reglas de interacción}

Se define el entorno de la partícula $i$ como el conjunto
$\mathcal{N}_i = \{\, j \neq i : d(\mathbf{r}_i, \mathbf{r}_j) < r_c \,\}$, donde
$d$ es la distancia mínima imagen bajo las condiciones periódicas y $r_c$ el
radio de interacción.

En el \textbf{modelo estándar}, $\Theta_i$ es el promedio vectorial de las
direcciones del entorno \textbf{incluyendo la de la propia partícula}, criterio
fijado por la cátedra:

\begin{equation}
    \Theta_i^{\,\text{Vicsek}}(t) = \arctan\!
    \frac{\sin\theta_i(t) + \sum_{j \in \mathcal{N}_i} \sin\theta_j(t)}
         {\cos\theta_i(t) + \sum_{j \in \mathcal{N}_i} \cos\theta_j(t)},
    \label{eq:vicsek}
\end{equation}

donde $\theta_i$ es la dirección de la partícula considerada y $\theta_j$ la de
cada uno de sus vecinos. El arco tangente se evalúa con la función de dos
argumentos, de modo de resolver el cuadrante correctamente. Una partícula sin
vecinos conserva su dirección más el ruido.

En el \textbf{modelo de votante}, la partícula no promedia: elige un único
elemento al azar y copia su dirección tal cual \cite{loscar2021},

\begin{equation}
    \Theta_i^{\,\text{votante}}(t) = \theta_k(t),
    \qquad
    k \sim \mathcal{U}\!\left(\mathcal{N}_i \cup \{i\}\right),
    \label{eq:votante}
\end{equation}

donde $k$ es un índice sorteado con probabilidad uniforme. Por simetría con la
Ec.~\eqref{eq:vicsek}, en la que la propia partícula participa del promedio, se
la incluye también en el sorteo de la Ec.~\eqref{eq:votante}; en consecuencia,
una partícula aislada conserva su dirección. El ángulo copiado es el del estado
en $t$: como las direcciones nuevas se escriben en un buffer separado y recién
se intercambian al terminar el paso, todas las copias de un mismo paso leen el
mismo estado.

\subsection{Observables}

La \textbf{polarización} es el módulo de la velocidad media normalizado por la
rapidez:

\begin{equation}
    \va(t) = \frac{1}{N v} \left\| \sum_{i=1}^{N} \mathbf{v}_i(t) \right\|
           = \frac{1}{N} \sqrt{ \Big(\sum_i \cos\theta_i\Big)^{\!2}
                              + \Big(\sum_i \sin\theta_i\Big)^{\!2} },
    \label{eq:polarizacion}
\end{equation}

donde $N$ es el número de partículas, $v$ la rapidez común y $\mathbf{v}_i$ la
velocidad de la partícula $i$. Vale $\va = 1$ si todas se mueven en la misma
dirección y $\va \approx 1/\sqrt{N}$ si las direcciones son independientes, que
es el piso por tamaño finito y no cero.

Un \textbf{cluster} es un conjunto de partículas tal que todo par está unido por
una cadena de saltos entre vecinos a distancia menor que $r_c$; es decir, una
componente conexa del grafo cuyos nodos son las partículas y cuyas aristas son
los pares dentro del radio de interacción. Llamando $\mathcal{C}_{\max}(t)$ a la
componente de mayor tamaño, el segundo observable es

\begin{equation}
    S(t) = \frac{|\mathcal{C}_{\max}(t)|}{N},
    \label{eq:componente}
\end{equation}

donde $|\mathcal{C}_{\max}|$ es la cantidad de partículas de la componente
gigante y $N$ el total. Se cumple $1/N \le S \le 1$.

% =========================================================
\section{Implementación}
% =========================================================

El trabajo se divide en dos módulos independientes comunicados por archivos de
texto, tal como exige la consigna: el motor de simulación, escrito en C++, y el
módulo de análisis y animación, escrito en Python. El motor nunca dibuja y el
módulo de animación nunca simula, de modo que la velocidad de reproducción de
las animaciones no queda supeditada a la velocidad de la simulación.

\subsection{Motor de simulación}

Cada paso de simulación consta de tres etapas:

\begin{enumerate}
    \item \textbf{Búsqueda de vecinos.} Se construye la lista $\mathcal{N}_i$
          para todas las partículas mediante el CIM.
    \item \textbf{Cálculo de las direcciones nuevas.} Se evalúa
          $\Theta_i$ según la Ec.~\eqref{eq:vicsek} o la
          Ec.~\eqref{eq:votante} y se le suma el ruido de la
          Ec.~\eqref{eq:angulo}. El resultado se almacena en un arreglo
          auxiliar, sin pisar el estado en $t$.
    \item \textbf{Desplazamiento.} Se aplica la Ec.~\eqref{eq:posicion} usando
          las direcciones \textbf{viejas} y se reingresan las coordenadas a la
          caja. Recién entonces el arreglo auxiliar pasa a ser el estado
          vigente.
\end{enumerate}

El uso del arreglo auxiliar es lo que garantiza que la actualización sea
efectivamente sincrónica: sin él, las partículas de índice alto se alinearían
con direcciones ya actualizadas de las de índice bajo, introduciendo una
dependencia espuria con el orden de recorrido.

\subsection{Búsqueda de vecinos}

Se reutiliza sin modificaciones el CIM implementado en el Trabajo Práctico 1,
junto con la geometría periódica y el parseo de argumentos, de manera que los
tiempos de ambos trabajos midan el mismo algoritmo. La caja se subdivide en
$M \times M$ celdas y cada partícula se compara únicamente contra las de su
propia celda y las ocho adyacentes. Para que ese vecindario sea suficiente debe
cumplirse $L/M > r_c$, lo que para $L = 10$ y $r_c = 1$ acota $M \le 9$; se
emplea $M = 9$, el máximo admisible, que es el que minimiza la cantidad de
partículas por celda. Cada par se registra una única vez, recorriendo sólo las
celdas vecinas de índice superior, y las distancias se calculan con el criterio
de mínima imagen.

Como control de correctitud, el motor implementa además una búsqueda por fuerza
bruta $\mathcal{O}(N^2)$: con la misma semilla, ambos métodos producen
trayectorias idénticas byte a byte, lo que verifica que el CIM no pierde ni
duplica pares.

\subsection{Formato de salida}

El motor produce archivos de texto plano:

\begin{itemize}
    \item \texttt{static.txt}: parámetros de la corrida, un par
          \texttt{clave valor} por línea.
    \item \texttt{dynamic.txt}: por cada cuadro, una línea con $t$ seguida de
          $N$ líneas \texttt{x y theta}.
    \item \texttt{observables.txt}: una línea \texttt{t va S} por paso.
\end{itemize}

Las series de observables se escriben en todos los pasos, mientras que los
cuadros de posiciones pueden ralearse: con $\rho = 8$ y \num{20000} pasos,
guardar todos los cuadros ocupa unos \SI{530}{\mega\byte}, frente a las
\num{20001} líneas de la serie de observables. Por eso el barrido de parámetros
guarda únicamente observables, y las animaciones se generan a partir de corridas
puntuales.

\subsection{Determinación del estado estacionario}

Cada corrida parte de una configuración desordenada (posiciones y ángulos
uniformes e independientes) y atraviesa un transitorio antes de alcanzar el
estado estacionario. Como el valor escalar de cada observable debe medirse
únicamente en el estacionario, se emplea un criterio automático y uniforme para
todas las corridas: se toma como referencia el promedio de $\va$ sobre la
segunda mitad de la serie --- que ya pertenece al estacionario --- y se define
$t_{eq}$ como el primer instante en que la serie suavizada por bloques cruza esa
referencia. Al ser el transitorio monótono, el cruce marca su final sin
depender de umbrales arbitrarios.

El mismo $t_{eq}$ se aplica a $\va$ y a $S$, de modo que ambos escalares
provienen de la misma ventana temporal. Los promedios se toman desde $t_{eq}$
hasta el final de la corrida.

% =========================================================
\section{Simulaciones}
% =========================================================

La Tabla~\ref{tab:parametros} resume los parámetros empleados; las densidades
se traducen en el número de partículas mediante la Ec.~\eqref{eq:densidad}. Todos los
resultados de las Secciones \ref{sec:temporal} a \ref{sec:vaS} provienen del
mismo barrido, de modo que las evoluciones temporales mostradas son corridas
efectivamente incluidas en las curvas escalares.

\begin{table}[H]
    \centering
    \caption{Parámetros de las simulaciones. Las semillas son consecutivas y se
    reutilizan entre modelos: para un mismo $\rho$ y una misma semilla, el
    modelo estándar y el de votante parten de la \textbf{misma} condición
    inicial, lo que hace que la comparación del ítem (f) sea pareada.}
    \label{tab:parametros}
    \begin{tabular}{l l}
        \toprule
        \textbf{Parámetro} & \textbf{Valor} \\
        \midrule
        Lado de la caja $L$                & \num{10} \\
        Contorno                           & periódico \\
        Radio de interacción $r_c$         & \num{1} \\
        Rapidez $v$                        & \num{0.03} \\
        Paso temporal $\Delta t$           & \num{1} \\
        Densidades $\rho$                  & \num{2}, \num{4}, \num{8} \\
        Partículas $N$                     & \num{200}, \num{400}, \num{800} \\
        Ruido $\eta$                       & de \num{0} a \num{5}, paso \num{0.5} (11 valores) \\
        Modelos                            & estándar, votante \\
        Pasos por corrida                  & \num{5000} \\
        Realizaciones $M$ por punto        & \num{10} (semillas 1 a 10) \\
        Celdas del CIM $M_{\text{CIM}}$    & \num{9} (máximo admisible) \\
        Corridas totales                   & \num{660} \\
        \bottomrule
    \end{tabular}
\end{table}

Cada punto de las curvas escalares se obtiene promediando \textbf{primero en el
tiempo} dentro del estacionario de cada realización y \textbf{luego entre las
$M = 10$ realizaciones}. La barra de error reportada en todas las figuras es el
\textbf{desvío estándar entre realizaciones}, no el error estándar de la media,
por lo que mide la dispersión efectiva entre condiciones iniciales y no la
precisión del promedio.

% =========================================================
\section{Resultados}
% =========================================================

\subsection{Animaciones características}

Las animaciones representan cada partícula mediante un vector con origen en su
posición y coloreado según el ángulo de su velocidad, usando un mapa de color
cíclico para que direcciones próximas a $-\pi$ y a $\pi$ reciban colores
próximos. La longitud del vector es fija: con $v = 0.03$ y $L = 10$ el
desplazamiento real por paso es el \SI{0.3}{\percent} del ancho de la caja y las
flechas se verían como puntos. Como la rapidez es común a todas las partículas,
dibujarlas con longitud uniforme no oculta información: la flecha indica
dirección, no módulo. La Tabla~\ref{tab:animaciones} reúne los casos
animados y sus enlaces.

\begin{table}[H]
    \centering
    \caption{Animaciones características. Todas corresponden al modelo estándar
    salvo donde se indica, con $L = 10$ y contorno periódico.}
    \label{tab:animaciones}
    \begin{tabular}{l l l}
        \toprule
        \textbf{Caso} & \textbf{Parámetros} & \textbf{Enlace} \\
        \midrule
        Ruido bajo (ordenado)      & $\rho = 4$, $\eta = 0.5$ & \texttt{[PENDIENTE]} \\
        Ruido alto (desordenado)   & $\rho = 4$, $\eta = 3$   & \texttt{[PENDIENTE]} \\
        Votante, ruido bajo        & $\rho = 4$, $\eta = 1$   & \texttt{[PENDIENTE]} \\
        \bottomrule
    \end{tabular}
\end{table}

% TODO: reemplazar [PENDIENTE] por los enlaces de YouTube (no subir como shorts).

Con $\eta$ bajo el sistema converge a una única bandada que recorre la caja con
una dirección común, reflejada en un color prácticamente uniforme. Con $\eta$
alto los colores se mezclan y no se distingue ninguna dirección preferencial,
aunque las partículas siguen formando grupos transitorios que se arman y se
deshacen. La Fig.~\ref{fig:cuadros} muestra esa evolución en cuadros fijos, que
es lo que puede imprimirse en este documento.

\begin{figure}[H]
    \centering
    \includegraphics[width=\textwidth]{cuadros_vicsek.png}
    \caption{Cuadros del modelo estándar a distintos instantes, partiendo de una
    configuración desordenada. Cada flecha es una partícula, con origen en su
    posición y color según el ángulo de su velocidad (barra de color a la
    derecha). Los parámetros de la corrida se indican en el encabezado de la
    figura.}
    \label{fig:cuadros}
\end{figure}

\subsection{Evolución temporal de los observables}
\label{sec:temporal}

La Fig.~\ref{fig:temporal_eta} muestra la evolución temporal de $\va$ y de $S$
para el modelo estándar a $\rho = 4$ y tres valores de $\eta$. En todos los
casos la polarización parte del valor correspondiente a una configuración
desordenada, crece durante un transitorio y luego fluctúa alrededor de un valor
constante. La línea vertical punteada marca el $t_{eq}$ determinado por el
criterio de la Sección 3.4 y la línea horizontal rayada el promedio temporal
tomado a partir de ese instante. Se aprecia que el transitorio es más largo
cuanto menor es el ruido: para $\eta = 0$ el sistema tarda del orden de
\num{700} pasos en ordenarse por completo, mientras que para $\eta \ge 4$ el
estacionario se alcanza en menos de \num{200} pasos.

\begin{figure}[H]
    \centering
    \includegraphics[width=\textwidth]{temporal_eta_vicsek_rho4.png}
    \caption{Evolución temporal de $\va$ (panel superior) y de $S$ (panel
    inferior) para el modelo estándar con $\rho = 4$ ($N = 400$), $L = 10$,
    $r_c = 1$, $v = 0.03$, semilla 1, para tres valores de $\eta$. La línea
    vertical punteada indica el inicio del estado estacionario $t_{eq}$ y la
    horizontal rayada el promedio temporal a partir de $t_{eq}$. El mismo
    $t_{eq}$, determinado sobre $\va$, se aplica a $S$.}
    \label{fig:temporal_eta}
\end{figure}

El panel inferior de la Fig.~\ref{fig:temporal_eta} muestra que $S$ alcanza
valores próximos a la unidad casi de inmediato y se mantiene allí durante toda
la corrida, con fluctuaciones mucho menores que las de $\va$. Esto no es un
artefacto: con $r_c = 1$, el número medio de vecinos por partícula es
$\langle k \rangle = \rho \pi r_c^2$, que vale \num{6.3}, \num{12.6} y \num{25.1}
para $\rho = 2$, \num{4} y \num{8} respectivamente. Todos esos valores están muy
por encima del umbral de percolación de un grafo geométrico aleatorio en dos
dimensiones ($\langle k \rangle \approx 4.5$), de modo que el sistema está
percolado desde la condición inicial y para cualquier $\eta$. Esta observación
condiciona la lectura de los ítems (d) y (e) y se retoma en la
Sección~\ref{sec:vaS}.

La Fig.~\ref{fig:temporal_modelos} compara las evoluciones de ambos modelos
desde la misma condición inicial. La diferencia es inmediata: el modelo estándar
alcanza $\va \approx 0.79$ mientras que el de votante cae a $\va \approx 0.08$
para los mismos $\rho$ y $\eta$, sin transitorio de ordenamiento apreciable.

\begin{figure}[H]
    \centering
    \includegraphics[width=\textwidth]{temporal_modelos_rho4_eta2.png}
    \caption{Comparación de la evolución temporal de $\va$ y $S$ entre el modelo
    estándar y el de votante, con $\rho = 4$ ($N = 400$), $\eta = 2$, $L = 10$,
    $r_c = 1$, $v = 0.03$ y la misma semilla, de modo que ambas corridas parten
    de idéntica condición inicial. Líneas verticales punteadas: $t_{eq}$ de cada
    modelo; horizontales rayadas: promedios en el estacionario.}
    \label{fig:temporal_modelos}
\end{figure}

\subsection{Polarización en función del ruido}
\label{sec:va}

La Fig.~\ref{fig:va_eta} presenta el observable escalar $\va$, definido en la
Ec.~\eqref{eq:polarizacion}, en función de $\eta$, para las tres densidades y
los dos modelos.

\begin{figure}[H]
    \centering
    \includegraphics[width=\textwidth]{va_vs_eta.png}
    \caption{Polarización $\va$ en el estado estacionario en función del ruido
    $\eta$, para $\rho = 2$, \num{4} y \num{8} y los dos modelos. $L = 10$,
    $r_c = 1$, $v = 0.03$, $\Delta t = 1$, \num{5000} pasos por corrida.
    Cada punto es el promedio temporal en el estacionario promediado sobre
    $M = 10$ realizaciones; la barra de error es el desvío estándar entre
    realizaciones. Líneas llenas y marcadores llenos: modelo estándar; líneas
    de trazos y marcadores huecos: modelo de votante.}
    \label{fig:va_eta}
\end{figure}

En el \textbf{modelo estándar} la polarización decae de manera continua desde
$\va = 1$ en $\eta = 0$ hasta valores próximos al piso por tamaño finito en
$\eta = 5$, sin salto abrupto, lo que es compatible con la transición continua
descripta en \cite{vicsek1995}. El efecto de la densidad es sistemático: a
igual ruido, mayor densidad implica mayor polarización. Para $\eta = 3$, por
ejemplo, $\va$ pasa de \num{0.47(4)} con $\rho = 2$ a \num{0.562(17)} con
$\rho = 4$ y \num{0.615(4)} con $\rho = 8$. La interpretación es directa a
partir de la Ec.~\eqref{eq:vicsek}: promediar sobre
$\langle k \rangle \propto \rho$ vecinos atenúa el ruido en un factor
$\sim 1/\sqrt{\langle k \rangle}$, de modo que aumentar la densidad equivale a
reducir el ruido efectivo. Consistentemente, la caída de la curva se desplaza
hacia $\eta$ mayores al aumentar $\rho$.

El \textbf{modelo de votante} se comporta de manera cualitativamente distinta.
La polarización se desploma apenas se enciende el ruido: basta $\eta = 0.5$
para que $\va$ caiga de \num{1} a \num{0.382(21)} con $\rho = 2$, y ya en
$\eta \approx 1.5$ el sistema está esencialmente desordenado. Más notable aún,
\textbf{el efecto de la densidad se invierte}: a igual $\eta$, mayor densidad
produce \textit{menor} polarización, exactamente al revés que en el modelo
estándar. Para $\eta = 1$, $\va$ vale \num{0.201(5)} con $\rho = 2$,
\num{0.151(4)} con $\rho = 4$ y \num{0.112(5)} con $\rho = 8$.

Ambos hechos se explican por el mismo mecanismo. Copiar a un único vecino no
promedia nada, de modo que el ruido de la Ec.~\eqref{eq:angulo} se incorpora
íntegro en cada paso, sin la atenuación $1/\sqrt{\langle k \rangle}$ que
proporciona la Ec.~\eqref{eq:vicsek}; aumentar el número de vecinos no mejora la
estimación de la dirección local, sólo acelera la propagación del desorden a
través del sistema. En $\eta = 5$ la polarización del votante alcanza
\num{0.0643}, \num{0.0454} y \num{0.0323} para $\rho = 2$, \num{4} y \num{8},
valores que coinciden dentro del \SI{10}{\percent} con el piso por tamaño finito
$1/\sqrt{N} = $ \num{0.0707}, \num{0.05} y \num{0.0354}: el sistema está
completamente desordenado. El modelo estándar, en cambio, se mantiene por encima
de ese piso en todo el rango explorado (\num{0.0748} frente a \num{0.0354} para
$\rho = 8$ en $\eta = 5$), es decir que conserva orden residual incluso en el
extremo del barrido.

\subsection{Componente gigante}
\label{sec:S}

La Fig.~\ref{fig:temporal_rho} muestra la evolución temporal de la fracción $S$
definida en la Ec.~\eqref{eq:componente}, para las tres densidades, y la Fig.~\ref{fig:S_eta} el valor escalar correspondiente en
función de $\eta$, obtenido con el mismo procedimiento que la
Fig.~\ref{fig:va_eta}.

\begin{figure}[H]
    \centering
    \includegraphics[width=\textwidth]{temporal_rho_vicsek_eta2.png}
    \caption{Evolución temporal de $\va$ y $S$ para el modelo estándar con
    $\eta = 2$ y las tres densidades $\rho = 2$, \num{4} y \num{8}
    ($N = 200$, \num{400} y \num{800}), $L = 10$, $r_c = 1$, semilla 1. Líneas
    verticales punteadas: $t_{eq}$; horizontales rayadas: promedio en el
    estacionario.}
    \label{fig:temporal_rho}
\end{figure}

\begin{figure}[H]
    \centering
    \includegraphics[width=\textwidth]{S_vs_eta.png}
    \caption{Fracción $S$ de partículas en la componente gigante, en el estado
    estacionario, en función de $\eta$, para las tres densidades y los dos
    modelos. Mismos parámetros, criterio de promediado y convención de barras de
    error que la Fig.~\ref{fig:va_eta}. Nótese que la escala vertical cubre todo
    el rango $[0,1]$ y que los datos ocupan únicamente su parte superior.}
    \label{fig:S_eta}
\end{figure}

El resultado central es que $S \approx 1$ en prácticamente todo el dominio
explorado. El mínimo sobre todo el barrido del modelo estándar es
$S = \num{0.914(13)}$, en $\rho = 2$ y $\eta = 3.5$; para $\rho = 8$ el valor
nunca baja de \num{0.999}. Como se argumentó en la Sección~\ref{sec:temporal},
esto es consecuencia directa de los parámetros fijados por la consigna: con
$L = 10$ y $\rho \ge 2$ el grafo de vecinos está muy por encima del umbral de
percolación, de modo que la componente gigante abarca casi todo el sistema
independientemente del ruido y del grado de alineación.

Se observa, no obstante, una estructura no monótona apreciable a $\rho = 2$, que
es la densidad más cercana al umbral y por lo tanto la única con margen para
mostrarla: $S$ decrece desde \num{1} hasta el mínimo de \num{0.914} en torno a
$\eta \approx 3.5$ y luego vuelve a crecer hasta \num{0.984} en $\eta = 5$. El
mínimo coincide aproximadamente con la región donde $\va$ cae más rápido. La
explicación es que en el régimen ordenado las partículas viajan juntas y
mantienen la conectividad, y en el régimen fuertemente desordenado se
distribuyen de manera prácticamente uniforme por la caja, lo que también
garantiza la percolación a esta densidad; es en la zona intermedia donde el
sistema forma grupos densos separados por regiones vacías, y por lo tanto donde
la componente gigante pierde algunos nodos.

\subsection{Relación entre polarización y componente gigante}
\label{sec:vaS}

La Fig.~\ref{fig:va_S} grafica $\va$ en función de $S$, con un punto por cada
valor de $\eta$.

\begin{figure}[H]
    \centering
    \includegraphics[width=\textwidth]{va_vs_S.png}
    \caption{Polarización $\va$ en función de la fracción $S$ de partículas en
    la componente gigante, en el estado estacionario. Cada punto corresponde a
    un valor de $\eta$ del barrido; las barras horizontales y verticales son los
    desvíos estándar entre las $M = 10$ realizaciones de $S$ y de $\va$
    respectivamente. Se distinguen las tres densidades por color y los dos
    modelos por estilo de línea y marcador.}
    \label{fig:va_S}
\end{figure}

La figura debe leerse teniendo presente la limitación establecida en la
Sección~\ref{sec:S}: como $S$ está confinado al intervalo $[0.91, 1]$ mientras
que $\va$ recorre casi todo $[0,1]$, los puntos se acumulan contra el borde
derecho y la relación entre ambos observables no puede caracterizarse en este
rango de parámetros. Dicho de otro modo, \textbf{con $L = 10$ y
$\rho = 2,\,4,\,8$ la conectividad espacial no es un grado de libertad efectivo}:
el sistema está siempre percolado y el ordenamiento angular ocurre sin que la
componente gigante cambie apreciablemente. La conclusión cuantitativa que sí
puede extraerse es justamente ésa: en el régimen estudiado, alineación y
conectividad están desacopladas, y $\va$ no es una función de $S$.

Un régimen en el que ambos observables sí varían de forma acoplada requiere
acercarse al umbral de percolación, lo que exige salir de los parámetros de la
consigna. A modo de verificación complementaria se corrió el modelo estándar con
$L = 40$ y $\rho = 0.6$ ($\langle k \rangle \approx 1.9$, por debajo del
umbral), obteniéndose un estado estacionario con $\va \approx 0.98$ y
$S \approx 0.85$ sostenidos: una bandada dominante acompañada de grupos satélite
espacialmente separados pero alineados con ella. Este caso confirma que la
degeneración observada es propia de los parámetros fijados y no del modelo ni de
la implementación, pero por estar fuera de la consigna no se incorpora a las
curvas del barrido principal.

\subsection{Comparación entre modelos}

Las Figuras \ref{fig:temporal_modelos}, \ref{fig:va_eta}, \ref{fig:S_eta} y
\ref{fig:va_S} incluyen ambos modelos superpuestos, según pide el ítem (f). Los
resultados se resumen así:

\begin{itemize}
    \item \textbf{En $\eta = 0$ ambos modelos son equivalentes}: los dos
          alcanzan $\va = 1$. Sin ruido, copiar a un vecino ya alineado y
          promediar vecinos ya alineados producen el mismo resultado, lo que
          sirve además como verificación cruzada de las dos implementaciones.
    \item \textbf{Para $\eta > 0$ las curvas se separan de inmediato.} El modelo
          de votante es dramáticamente más sensible al ruido: en $\eta = 1$ y
          $\rho = 4$, el estándar conserva $\va = \num{0.945(2)}$ mientras que
          el votante ya cayó a $\va = \num{0.151(4)}$.
    \item \textbf{La dependencia con la densidad tiene signo opuesto}, como se
          detalló en la Sección~\ref{sec:va}: creciente para el estándar,
          decreciente para el votante.
    \item \textbf{En $S$ la diferencia es menor}, porque ambos modelos operan
          muy por encima del umbral de percolación. El votante presenta valores
          de $S$ ligeramente inferiores a $\rho = 2$ ($S = \num{0.941(20)}$ en
          $\eta = 1$, frente a \num{0.988(5)} del estándar), consistente con que
          su desorden angular reparte las partículas más uniformemente.
\end{itemize}

\subsection{Tiempos de ejecución del CIM}

Para el ítem (g) se midió el tiempo por búsqueda de vecinos del CIM y se lo
comparó con el registrado en el Trabajo Práctico 1. El cronómetro abarca
únicamente la búsqueda --- incluido el armado de las celdas --- y excluye la
integración de la dinámica y la escritura a disco.

La comparación exige igualar la configuración geométrica: las mediciones del
Trabajo Práctico 1 se realizaron con $L = 20$ y $M = 13$, mientras que el
estudio del presente trabajo emplea $L = 10$ y $M = 9$. A igual $N$ ambas
configuraciones corresponden a densidades distintas, de modo que los tiempos no
son directamente comparables. Por eso las corridas de esta sección se ejecutan
con $L = 20$ y $M = 13$, replicando la caja del Trabajo Práctico 1, sobre el
mismo conjunto de valores de $N$ (de \num{10} a \num{1150}).

\textbf{Nota metodológica.} Se verificó que las mediciones de tiempo son
extremadamente sensibles a la carga de la máquina. Al reejecutar el binario
\textit{sin modificar} del Trabajo Práctico 1 sobre un equipo con otros procesos
en ejecución, el tiempo para $N = 1150$ pasó de los \SI{0.904}{\milli\second}
originales a \SI{3.756}{\milli\second}, con un desvío del \SI{56}{\percent} de
la media. En consecuencia, ambos conjuntos de mediciones deben tomarse con el
equipo descargado para que la comparación tenga sentido.

% TODO: reejecutar con la maquina descargada y actualizar la figura:
%   cd engine
%   ./build/flock bench --l 20 --ms 13 \
%       --ns 10,113,217,320,424,528,631,735,839,942,1046,1150 \
%       --steps 200 --repeats 5 --eta 6.283185307 --out ../data/bench_l20.txt
%   python3 visualization/bench.py --bench data/bench_l20.txt --tp1 data/tp1.txt --log

\begin{figure}[H]
    \centering
    \includegraphics[width=\textwidth]{tiempos_cim.png}
    \caption{Tiempo por búsqueda de vecinos en función de $N$, con $L = 20$,
    $r_c = 1$ y $M = 13$. Se comparan el CIM y la fuerza bruta del presente
    trabajo con el CIM del Trabajo Práctico 1. Las barras de error son el desvío
    estándar entre repeticiones.}
    \label{fig:tiempos}
\end{figure}

Independientemente del valor absoluto, la comparación entre el CIM y la fuerza
bruta dentro de una misma sesión de medición --- y por lo tanto bajo idénticas
condiciones de carga --- confirma el comportamiento esperado: el costo de la
fuerza bruta crece como $\mathcal{O}(N^2)$ mientras que el del CIM crece de
manera aproximadamente lineal con $N$ a densidad fija, ya que el número medio de
partículas por celda permanece constante y cada partícula se compara sólo contra
un entorno acotado.

% =========================================================
\section{Conclusiones}
% =========================================================

A partir de los resultados presentados en la Sección 5 se concluye:

\begin{enumerate}
    \item \textbf{El modelo estándar exhibe una transición continua de orden a
          desorden controlada por $\eta$} (Fig.~\ref{fig:va_eta}). La
          polarización decae suavemente desde $\va = 1$ en $\eta = 0$ y se
          mantiene por encima del piso por tamaño finito $1/\sqrt{N}$ en todo el
          rango explorado, conservando orden residual incluso en $\eta = 5$.

    \item \textbf{En el modelo estándar la densidad favorece el orden.} A
          $\eta = 3$, $\va$ crece de \num{0.47(4)} a \num{0.615(4)} al pasar de
          $\rho = 2$ a $\rho = 8$ (Fig.~\ref{fig:va_eta}), y la caída de la
          curva se desplaza hacia ruidos mayores. Esto es consistente con que
          el promedio de la Ec.~\eqref{eq:vicsek} atenúa el ruido en un factor
          $\sim 1/\sqrt{\rho \pi r_c^2}$.

    \item \textbf{El modelo de votante es mucho más sensible al ruido y su
          dependencia con la densidad tiene signo opuesto.} Con $\rho = 4$ y
          $\eta = 1$, el votante alcanza $\va = \num{0.151(4)}$ frente a
          $\va = \num{0.945(2)}$ del estándar; y a $\eta = 1$ la polarización
          del votante \textit{disminuye} de \num{0.201(5)} a \num{0.112(5)} al
          aumentar $\rho$ de \num{2} a \num{8} (Fig.~\ref{fig:va_eta}). Ambos
          hechos responden a que copiar un único vecino no promedia el ruido,
          por lo que más vecinos no mejoran la estimación de la dirección local
          y sólo aceleran la propagación del desorden. En $\eta = 5$ el votante
          alcanza el piso $1/\sqrt{N}$ dentro del \SI{10}{\percent}, es decir
          desorden completo.

    \item \textbf{Ambos modelos coinciden exactamente en $\eta = 0$}
          ($\va = 1$), lo que verifica de forma cruzada las dos
          implementaciones.

    \item \textbf{En el rango de parámetros de la consigna la componente gigante
          no es un grado de libertad efectivo.} Con $L = 10$ y $\rho \ge 2$ se
          tiene $\langle k \rangle \ge 6.3$, muy por encima del umbral de
          percolación bidimensional, y en consecuencia $S \ge 0.91$ en todo el
          barrido (Fig.~\ref{fig:S_eta}). La polarización varía en casi todo
          $[0,1]$ mientras $S$ permanece confinado cerca de la unidad, de modo
          que \textbf{alineación y conectividad están desacopladas} y $\va$ no
          es función de $S$ (Fig.~\ref{fig:va_S}). La única estructura
          apreciable es un mínimo de $S \approx \num{0.914}$ en torno a
          $\eta \approx 3.5$ para $\rho = 2$, la densidad más próxima al umbral.

    \item \textbf{El CIM reduce el costo de la búsqueda de vecinos respecto de
          la fuerza bruta}, con un crecimiento aproximadamente lineal en $N$ a
          densidad fija frente al $\mathcal{O}(N^2)$ de la fuerza bruta
          (Fig.~\ref{fig:tiempos}). La comparación con el Trabajo Práctico 1
          exige igualar tanto la geometría ($L$ y $M$) como la carga de la
          máquina: se comprobó que el mismo binario del Trabajo Práctico 1
          arroja tiempos hasta \num{4.2} veces mayores cuando el equipo no está
          descargado.
\end{enumerate}

% =========================================================
\section*{Referencias}
\addcontentsline{toc}{section}{Referencias}
% =========================================================

\begin{thebibliography}{9}

\bibitem{vicsek1995}
T. Vicsek, A. Czirók, E. Ben-Jacob, I. Cohen y O. Shochet,
``Novel type of phase transition in a system of self-driven particles'',
Physical Review Letters,
vol. 75, nro. 6, p. 1226 (1995).

\bibitem{loscar2021}
E. S. Loscar, G. Baglietto y F. Vazquez,
``Noisy multistate voter model for flocking in finite dimensions'',
Physical Review E,
vol. 104, nro. 3, p. 034111 (2021).

\end{thebibliography}

\end{document}
